\cleardoublepage
\thispagestyle{empty}
\bigtitle{S.\@ Antonii de Padua Confessoris et Eccl.\@ Doct.}
\addcontentsline{toc}{section}{XIII Junii} \phantomsection
\duplex

\header{xiii junii 2026}

\office{ad missam}
%\label{missa_doct}
\addcontentsline{toc}{subsection}{Ad Missam} \phantomsection
\proper{Introitus}{Eccli. 15, 5; Ps. 91, 2}

\texteparacol{\lettrine{I}{n}  médio Ecclésiæ apéruit os ejus: et implévit eum Dóminus spíritu sapiéntiæ et intelléctus: stolam glóriæ índuit eum.
℣. Bonum est confitéri Dómino: et psállere nómini tuo, Altíssime.
Glória Patri. Sicut erat.
In médio.}{english}{\noindent In the midst of the assembly he opened his mouth; and the Lord filled him with the spirit of wisdom and understanding; he clothed him with a robe of glory. ℣. It is good to give thanks to the Lord, to sing praise to thy name, Most High. Glory be to the Father. As it was in the beginning. In the midst.}

\masspart{Oratio}

\texteparacol{\lettrine{E}{c}clésiam tuam, Deus, beáti Antónii Confessóris atque Do\-ctóribus tui solé\-mnitas votíva lætíficet:~† ut spirituálibus semper muniátur auxíliis,~* et gáudiis pérfrui mereátur ætérnis. Per Dóminum.

\lettrine[ante=2]{D}{e}us, qui nobis in Corde Fílii tui, nostris vulneráto peccátis, infinítos dile\-ctiónis thesáuros misericórditer largíri dignáris:~† concéde quǽsumus; ut illi devótum pietátis nostræ præstántes obséquium,~* dignæ quoque satisfa\-ctiónis exhibeámus offícium. Per eúmdem Dóminum.}{english}{\noindent May the solemn feast of blessed Anthony, thy confessor and Doctor, make thy church rejoice, o God, so that, ever sustained by spiritual help, she may deserve to reap everlasting joy.
Through our Lord.

\noindent 2. O God, who hast suffered the heart of thy Son to be wounded by our sins, and in that very heart hast bestowed on us the abundant riches of thy love: grant that the devout homage of our hearts, which we render unto him, may of thy mercy be deemed a recompence acceptable in thy sight. Through the same Lord.}

\reading{Epistola}{2 Tim. 4, 1-- 8}

\texteparacol{\lettrine{C}{a}ríssime: Testíficor coram Deo, et Jesu Christo, qui judicatúrus est vivos et mórtuos, per advéntum i\-psíus et regnum ejus: prǽdica verbum, insta opportúne, importune: árgue, óbsecra, íncrepa in o\-mni patiéntia, et doctrína. Erit enim tempus, cum sanam doctrínam non sustinébunt, sed ad sua desidéria, coacervábunt sibi magistros, pruriéntes áuribus, et a veritáte quidem audítum avértent, ad fábulas autem converténtur. Tu vero vígila, in ó\-mnibus labóra, opus fac Evangelístæ, ministérium tuum imple. Sóbrius esto. Ego enim jam delíbor, et tempus resolutiónis meæ instat. Bonum certámen certávi, cursum consummávi, fidem servávi. In réliquo repósita est mihi coróna justítiæ, quam reddet mihi Dóminus in illa die, justus judex: non solum autem mihi, sed et iis, qui díligunt advéntum ejus.}{english}{\noindent Beloved: I charge you, in the sight of God and Christ Jesus, who will judge the living and the dead by his coming and by his kingdom, preach the word, be urgent in season, out of season; reprove, entreat, rebuke with all patience and teaching. For there will come a time when they will not endure the sound doctrine; but having itching ears, will heap up to themselves teachers according to their own lusts, and they will turn away their hearing from the truth and turn aside rather to fables. But be watchful in all things, bear with tribulation patiently, work as a preacher of the Gospel, fulfill your ministry. As for me, I am already being poured out in sacrifice, and the time of my deliverance is at hand. I have fought the good fight, I have finished the course, I have kept the faith. For the rest, there is laid up for me a crown of justice, which the Lord, the just Judge, will give to me in that day; yet not to me only, but also to those who love his coming.}

\proper{Graduale \& Alleluia}{Ps. 36, 30 -- 31; Eccli. 45, 9}

\texteparacol{\lettrine{O}{s} justi meditábitur sapiéntiam, et lingua ejus loquétur judícium.
℣. Lex Dei ejus in corde i\-psíus: et non supplantabúntur gressus ejus.

\lettrine{A}{l}elúia, allelúia. Amávit eum Dóminus, et ornávit eum: stolam glóriæ índuit eum. Allelúia.}{english}{\noindent The mouth of the just shall meditate wisdom and his tongue shall speak judgment.
℣. The law of his God is in his heart, and his steps shall not be supplanted.

\noindent Alleluia, alleluia. The Lord loved him and beautified him and clothed him with a robe of glory. Alleluia.}

\reading{Evangelium}{Matt. 5, 3 -- 19}

\texteparacol{\lettrine{I}{n} illo témpore: Dixit Jesus discípulis suis: Vos estis sal terræ. Quod si sal evanúerit, in quo saliétur? Ad níhilum valet ultra, nisi ut mittátur foras, et conculcétur ab homínibus. Vos estis lux mundi. Non potest cívitas abscóndi supra montem pósita. Neque accéndunt lucérnam, et ponunt eam sub módio, sed super candelábrum, ut lúceat ó\-mnibus, qui in domo sunt. Sic luceat lux vestra coram homínibus, ut vídeant ópera vestra bona, et gloríficent Patrem vestrum, qui in cælis est. Nolíte putáre, quóniam veni sólvere legem aut prophétas: non veni sólvere, sed adimplére. Amen, quippe dico vobis, donec tránseat cælum et terra, jota unum aut unus apex non præteríbit a lege, donec ó\-mnia fiant. Qui ergo sólverit unum de mandátis istis mínimis, et docúerit sic hómines, mínimus vocábitur in regno cælórum: qui autem fécerit et docúerit, hic magnus vocábitur in regno cælórum.}{english}{\noindent At that time, Jesus said to his disciples: You are the salt of the earth. But if the salt lose its savour, wherewith shall it be salted? It is good for nothing any more but to be cast out, and to be trodden on by men. You are the light of the world. A city seated on a mountain cannot be hid. Neither do men light a candle and put it under a bushel, but upon a candlestick, that it may shine to all that are in the house. So let your light shine before men, that they may see your good works, and glorify your Father who is in heaven. Do not think that I am come to destroy the law, or the prophets. I am not come to destroy, but to fulfill. For amen I say unto you, till heaven and earth pass, one jot, or one tittle shall not pass of the law, till all be fulfilled. He therefore that shall break one of these least commandments, and shall so teach men, shall be called the least in the kingdom of heaven. But he that shall do and teach, he shall be called great in the kingdom of heaven.}

\proper{Offertorium}{Ps. 91, 13}

\texteparacol{\lettrine{J}{u}stus ut palma florébit: sicut cedrus, quæ in Líbano est, multiplicábitur}{english}{\noindent The just shall flourish like the palm tree: he shall grow up like the cedar of Libanus.}

\masspart{Secreta}
\texteparacol{\lettrine{P}{r}æsens oblátio fiat, Dómine, pópulo tuo salutáris: pro quo dignátus es Patri tuo te vivéntem hóstiam immoláre:
Qui cum eódem Deo Patre et Spíritu San\-cto.

\lettrine[ante=2]{R}{e}spice, quǽsumus, Dómine, ad ineffábilem Cordis dilé\-cti Fílii tui caritátem: ut quod offérimus sit tibi munus accé\-ptum et nostrórum expiátio deli\-ctórum.
Per eúmdem Dóminum.}{english}{\noindent May the offering here before thee, o Lord, bring safety to thy people, for whom thou didst mercifully offer thyself a living victim to the Father,
Who, with the same God the Father and the Holy Spirit.

\noindent  2. Look, we beseech thee, o Lord, upon the heart of thy beloved Son, with its boundless love, so that what we offer, may be an acceptable gift in thy sight and an atonement for our sins.
Through the same Lord.}

\masspart{Præfatio de sacratissimo Cordis Jesu}

\texteparacol{\lettrine{V}{e}re dignum et justum est, æquum et salutáre, nos tibi semper et ubíque grátias ágere: Dómine san\-cte, Pater omnípotens, ætérne Deus: Qui Unigénitum tuum, in Cruce pendéntem, láncea mílitis transfígi voluísti: ut apértum Cor, divínæ largitátis sacrárium, torréntes nobis fúnderet miseratiónis et grátiæ: et, quod amóre nostri flagráre numquam déstitit, piis esset réquies et pæniténtibus pater et salútis refúgium. Et ídeo cum Angelis et Archángelis, cum Thronis et Dominatiónibus cumque omni milítia cæléstis exércitus hymnum glóriæ tuæ cánimus, sine fine dicéntes:}{english}{\noindent It is truly meet and just, right and for our salvation, that we should at all times, and in all places, give thanks unto Thee, O holy Lord, Father almighty, everlasting God; Who didst will that Thine only-begotten Son, while hanging on the cross, should be pierced by a soldier's spear, that the Heart thus opened, a shrine of divine bounty, should pour out on us streams of mercy and grace, and that what never ceased to burn with love for us, should be a resting-place to the devout, and open as a refuge of salvation to the penitent. And therefore with angels and archangels, with thrones and dominions, and with all the hosts of the heavenly army, we sing the hymn of thy glory, evermore saying:}

%\rubrique{Præfatio, \normaltext{\pageref{præfatio}.}}

\proper{Communio}{Luc. 12, 42}

\texteparacol{\lettrine{F}{i}délis servus et prudens, quem constítuit dóminus super famíliam suam: ut det illis in témpore trítici mensúram.}{english}{\noindent The faithful and prudent servant whom the master will set over his household to give them their ration of grain in due time.}

\newpage
\masspart{Postcommunio}
\texteparacol{\lettrine{D}{i}vínis, Dómine, munéribus satiáti: quǽsumus; ut, beáti Antónii Confessóris tui atque Do\-ctóris méritis et intercessióne, salutáris sacrifícii sentiámus effé\-ctum.
Per Dóminum.

\lettrine[ante=2]{P}{r}ǽbeant nobis, Dómine Jesu, divínum tua san\-cta fervórem: quo dulcíssimi Cordis tui suavitáte percé\-pta; discámus terréna despícere, et amáre cæléstia:
Qui vivis et regnas.}{english}{\noindent Refreshed by divine gifts, we beseech thee, o Lord, that through the merits and intercession of blessed Anthony, thy confessor and Doctor, we may experience the effect of the saving sacrifice.
Through our Lord.

\noindent 2. May thy sacrament, o Lord Jesus, give us holy zeal, so that, seeing the sweetness of thy most loving heart, we may learn to despise the things of earth and love those of heaven.
Who livest and reignest.}


\office{ad vesperas}
\addcontentsline{toc}{subsection}{Ad  Vesperas} \phantomsection

\gscore[]{℣.}{Deus_In_Adjutorium_Festivus}

\gscore[1. Ant.]{1. g}{suavi_jugo}
 %% watch in case initial remains low after a couple runs
\psalmus{109}{109_1g}{109_1}
 \smallscore{an_suavi_jugo}
 
 \gscore[2. Ant.]{2. D}{misericors_et_miserator}
\psalmus{110}{110_2}{110_2}
\smallscore{an_misericors_et_miserator}
 
  \gscore[3. Ant.]{7. a}{exortum_est_sacred_heart}
  \psalmus{111}{111_7a}{111_7}
  \smallscore{an_exortum_est_sacred_heart}
  
   \gscore[4 Ant.]{8. c}{quid_retribuam}
  \psalmus{115}{115_8c}{115_8}
  \smallscore{an_quid_retribuam}
   
    \gscore[5. Ant.]{4. A*}{apud_dominum_propitiatio}
    \psalmus{129}{129_4A_A_star}{129_4A_A_star}
    \smallscore{an_apud_dominum_propitiatio}

\capitulum{1 Petri 5, 6 – 7}

\lettrine{C}{a}ríssimi: Humiliámini sub poténti manu Dei, ut vos exáltet in témpore visitatiónis;~† omnem sollicitúdinem vestram proiciéntes in eum,~* quóniam ipsi cura est de vobis.

\rr Deo grátias.

\hymnus
   \gscore[]{3.}{hy_en_ut_superba_solesmes_1961}
   
   \smallscore{versiculus_dom_octave_sc}

  \admagnificat
  
  \gscore[]{1. f}{cognoverunt}
  
    \smallscore{Magnificat_1f}
    
  \magnificattexte{Magnificat_1_ordinaire}
  
  \smallscore{an_cognoverunt}

  \vv Dóminus vobíscum. \rr Et cum spíritu tuo.
  
  \oratio
  
  Orémus.
  
\lettrine{P}{r}oté\-ctor in te sperántium, Deus, sine quo nihil est válidum, nihil san\-ctum:~† multíplica super nos misericórdiam tuam; ut, te re\-ctóre, te duce, sic transeámus per bona temporália,~* ut non amittámus ætérna. Per Dóminum.

\rr Amen.

  \gscore[Ant.]{2.}{an_cs_o_doctor_optime_solesmes_1961_antoni}
  
  \vv Justum dedúxit Dóminus per vias rectas.
  
\rr Et osténdit illi regnum Dei.

\zerobaroffsettextleft
\zerobaroffsettextright
\gscore[Ant.]{1.}{an_cs_sacerdos_et_pontifex_ora_solesmes_1961}
\resetbaroffsettextleft
\resetbaroffsettextright

\vv Amávit eum Dóminus, et ornávit eum.

\rr Stolam glóriæ índuit eum.

Orémus.

\lettrine{E}{x}áudi, quǽsumus Dómine, preces nostras, quas in beáti Basilíi Confessóris tui atque Pontíficis sole\-mnitáte deférimus:~† et, qui tibi digne méruit famulári, ejus intercedéntibus méritis,~* ab ó\-mnibus nos absólve peccátis.

 \gscore[Ant.]{1.}{an_ad_jesum_autem_solesmes_1961}
 
 \vv Hauriétis aquas in gáudio.

\rr De fóntibus Salvatóris.

Orémus.

\lettrine{D}{e}us, qui nobis in Corde Fílii tui, nostris vulneráto peccátis, infinítos dile\-ctiónis thesáuros misericórditer largíri dignáris:~† concéde quǽsumus; ut illi devótum pietátis nostræ præstántes obséquium,~* dignæ quoque satisfa\-ctiónis exhibeámus offícium. Per eúmdem Dóminum.

\rr Amen.

   \vv Dóminus vobíscum. \rr Et cum spíritu tuo.
   
   \gscore[]{2.}{BD_Semiduplex_Vesp}
   
  \vv Fidélium ánimæ per misericórdiam Dei requiéscant in pace.
  
  \rr Amen.
  
 \rubrique{Tunc dicitur Completorium. Si non dicendum:}
  
  \secreto
  
\vv Dóminus det nobis suam pacem.
  
\rr Et vitam ætérnam. Amen.

\rubrique{\normaltext{Salve Regína,} etc., \normaltext{\pageref{an_salve_regina_solesmes}.}}

